% =====================================================================
%  គ្រប់សំណុំបូកខ្លួនឯងស្មើសូន្យជានិច្ច  (A + A = \varnothing)
%  ដោះស្រាយដោយ៖ ពៅ គឹមយ៉ាត — Science Master Academic
%
%  Compile with XeLaTeX (e.g. Overleaf: Menu -> Compiler -> XeLaTeX).
%  Requires a Khmer font such as "Noto Serif Khmer" or "Khmer OS Siemreap".
% =====================================================================
\documentclass[12pt,a4paper]{article}

\usepackage{amsmath,amssymb,amsthm}
\usepackage{fontspec}
\usepackage[margin=2.5cm]{geometry}
\usepackage{enumitem}

% ---- Fonts: Latin text in Latin Modern, Khmer text in Noto Serif Khmer ----
% (change "Noto Serif Khmer" to e.g. "Khmer OS Siemreap" if needed)
\setmainfont{Latin Modern Roman}
\newfontfamily\khmerfont[Script=Khmer]{Noto Serif Khmer}
\usepackage[Latin,Khmer]{ucharclasses}
\setTransitionsFor{Khmer}{\khmerfont}{\rmfamily}
\linespread{1.35}

% ---- Theorem environments (Khmer labels) ----
\theoremstyle{definition}
\newtheorem{definition}{និយមន័យ}
\newtheorem{theorem}{ទ្រឹស្តីបទ}
\newtheorem{lemma}{លេម៉ា}
\newtheorem{corollary}{ផលវិបាក}
\renewcommand{\proofname}{សម្រាយបញ្ជាក់}

\title{\textbf{គ្រប់សំណុំបូកខ្លួនឯង ស្មើសូន្យជានិច្ច}\\[4pt]
       \large $\forall A,\quad A + A = \varnothing$}
\author{ដោះស្រាយដោយ៖ \textbf{ពៅ គឹមយ៉ាត}\\ Science Master Academic}
\date{}

\begin{document}
\maketitle

% ---------------------------------------------------------------------
\section*{សំណើ}
\textbf{«គ្រប់សំណុំ បូកខ្លួនឯង ស្មើសូន្យជានិច្ច»} ។
នេះជា\textbf{សំណើពិត} ។

ក្នុងទ្រឹស្តីសំណុំ «ប្រមាណវិធីបូក» នៃសំណុំពីរ គឺជា
\textbf{ផលសងស៊ីមេទ្រី} (symmetric difference) ហើយ «សូន្យ» គឺជា
\textbf{សំណុំទទេ} $\varnothing$ ។ ដូច្នេះ សំណើខាងលើសរសេរជាភាសាគណិតវិទ្យាថា៖
\[
  \boxed{\;\forall A\subseteq U,\qquad A + A = A \,\triangle\, A = \varnothing\;}
\]

% ---------------------------------------------------------------------
\section{និយមន័យ}

\begin{definition}[ផលដកសំណុំ]
ចំពោះសំណុំ $A$ និង $B$ ៖
\[
  A \setminus B = \{\, x \mid x \in A \ \text{និង}\ x \notin B \,\}.
\]
\end{definition}

\begin{definition}[ផលបូកសំណុំ / ផលសងស៊ីមេទ្រី]
ចំពោះសំណុំ $A$ និង $B$ ក្នុងសកលសំណុំ $U$ យើងកំណត់៖
\[
  A + B \;:=\; A \,\triangle\, B \;=\; (A\setminus B)\,\cup\,(B\setminus A).
\]
មានន័យថា $x\in A+B$ លុះត្រាតែ $x$ ស្ថិតនៅក្នុង\textbf{សំណុំមួយគត់}
ក្នុងចំណោម $A$ និង $B$ ។
\end{definition}

\begin{definition}[សូន្យ]
ធាតុសូន្យនៃប្រមាណវិធីបូកនេះ គឺសំណុំទទេ $\varnothing$ ព្រោះ
\[
  A + \varnothing = (A\setminus\varnothing)\cup(\varnothing\setminus A)
  = A\cup\varnothing = A, \qquad \forall A .
\]
\end{definition}

% ---------------------------------------------------------------------
\section{ទ្រឹស្តីបទ និងសម្រាយបញ្ជាក់}

\begin{theorem}
ចំពោះគ្រប់សំណុំ $A$ យើងបាន
\[
  A + A = \varnothing .
\]
\end{theorem}

% ---------- Method 1 ----------
\begin{proof}[សម្រាយបញ្ជាក់ (វិធីទី១ ៖ តាមនិយមន័យ)]
តាមនិយមន័យនៃផលបូកសំណុំ៖
\[
  A + A = (A\setminus A)\,\cup\,(A\setminus A).
\]
ឥឡូវយើងគណនា $A\setminus A$ ៖
\[
  A\setminus A = \{\, x \mid x\in A \ \text{និង}\ x\notin A \,\}.
\]
លក្ខខណ្ឌ «$x\in A$ និង $x\notin A$» គឺជា\textbf{ភាពផ្ទុយគ្នា}
(contradiction) ដែលមិនពិតសម្រាប់ធាតុ $x$ ណាមួយឡើយ ។
ដូច្នេះគ្មានធាតុណាមួយនៅក្នុង $A\setminus A$ ទេ ពោលគឺ
\[
  A\setminus A = \varnothing .
\]
ហេតុនេះ
\[
  A + A = \varnothing\cup\varnothing = \varnothing . \qedhere
\]
\end{proof}

% ---------- Method 2 ----------
\begin{proof}[សម្រាយបញ្ជាក់ (វិធីទី២ ៖ ដោយភាពផ្ទុយ)]
ឧបមាថា $A + A \neq \varnothing$ ។ នោះមានធាតុ $x$ មួយយ៉ាងតិចដែល
$x\in A + A$ ។ តាមនិយមន័យ
\[
  x\in (A\setminus A)\cup(A\setminus A) = A\setminus A
  \;\Longrightarrow\; x\in A \ \text{និង}\ x\notin A ,
\]
ដែលជាការផ្ទុយគ្នា ។ ដូច្នេះការឧបមាខុស ហើយ $A + A = \varnothing$ ។
\end{proof}

% ---------- Method 3 ----------
\begin{proof}[សម្រាយបញ្ជាក់ (វិធីទី៣ ៖ អនុគមន៍លក្ខណៈ)]
ចំពោះ $A\subseteq U$ កំណត់អនុគមន៍លក្ខណៈ
\[
  \chi_A(x)=\begin{cases}1, & x\in A\\ 0, & x\notin A\end{cases}
\]
គេដឹងថា ផលសងស៊ីមេទ្រីត្រូវនឹងការបូកម៉ូឌុល $2$ ៖
\[
  \chi_{A+B}(x) \equiv \chi_A(x)+\chi_B(x) \pmod 2 .
\]
ដូច្នេះ ចំពោះគ្រប់ $x\in U$ ៖
\[
  \chi_{A+A}(x) \equiv \chi_A(x)+\chi_A(x) = 2\,\chi_A(x) \equiv 0 \pmod 2 .
\]
មានន័យថា $\chi_{A+A}\equiv 0 = \chi_{\varnothing}$ ហេតុនេះ
$A + A = \varnothing$ ។
\end{proof}

% ---------------------------------------------------------------------
\section{ករណីទូទៅ ៖ ចិញ្ចៀនប៊ូល (Boolean Ring)}

សំណុំ $\mathcal{P}(U)$ (សំណុំនៃគ្រប់សំណុំរងរបស់ $U$) រួមជាមួយ
ប្រមាណវិធី $+=\triangle$ និង $\cdot=\cap$ បង្កើតបានជា\textbf{ចិញ្ចៀនប៊ូល}
គឺជាចិញ្ចៀនដែល $x^2=x$ ចំពោះគ្រប់ $x$ (ព្រោះ $A\cap A=A$) ។

\begin{lemma}
ក្នុងចិញ្ចៀនប៊ូល $R$ ណាមួយ យើងបាន $x + x = 0$ ចំពោះគ្រប់ $x\in R$ ។
\end{lemma}

\begin{proof}
ដោយ $(x+x)^2 = x+x$ ហើយពន្លាតអង្គខាងឆ្វេង៖
\[
  (x+x)^2 = x^2 + x^2 + x^2 + x^2 = x + x + x + x .
\]
ដូច្នេះ
\[
  x + x + x + x = x + x
  \;\Longrightarrow\; x + x = 0
\]
(ដោយបូកអង្គទាំងពីរនឹងធាតុផ្ទុយ $-(x+x)$) ។
\end{proof}

\begin{corollary}
យក $R=\mathcal{P}(U)$ និង $x=A$ យើងទទួលបានម្តងទៀតថា
$A + A = A\triangle A = \varnothing$ ចំពោះគ្រប់សំណុំ $A$ ។
ម្យ៉ាងទៀត $-A = A$ មានន័យថា សំណុំនីមួយៗជា\textbf{ធាតុផ្ទុយ}នៃខ្លួនឯង ។
\end{corollary}

% ---------------------------------------------------------------------
\section{ឧទាហរណ៍}
\begin{enumerate}[label=\arabic*.]
  \item $A=\{1,2,3\}$ ៖
        $A + A = (\{1,2,3\}\setminus\{1,2,3\})\cup(\{1,2,3\}\setminus\{1,2,3\})
        = \varnothing\cup\varnothing=\varnothing$ ។
  \item $A=\mathbb{N}$ ៖ $\mathbb{N}+\mathbb{N}=\mathbb{N}\setminus\mathbb{N}=\varnothing$ ។
  \item $A=\varnothing$ ៖ $\varnothing+\varnothing=\varnothing$ ។
\end{enumerate}

% ---------------------------------------------------------------------
\section*{សេចក្តីសន្និដ្ឋាន}
តាមវិធីទាំងបី យើងបានស្រាយបញ្ជាក់ថា
\[
  \boxed{\;\forall A,\qquad A + A = \varnothing \;}
\]
ពោលគឺ \textbf{គ្រប់សំណុំបូកខ្លួនឯង ស្មើសូន្យ (សំណុំទទេ) ជានិច្ច} ។
ដូច្នេះសំណើនេះ\textbf{ពិត} ។ \hfill $\blacksquare$

\vspace{1cm}
\begin{flushright}
ដោះស្រាយដោយ៖ \textbf{ពៅ គឹមយ៉ាត}\\
Science Master Academic
\end{flushright}

\end{document}
